\section{任务}
\label{sec:tasks}

\wcabench{} 定义五个任务，构成从监督回归到因果推断的连续难度谱。每个任务都以相同的五要素刻画——形式化定义、输入与输出、指标、基线与领域挑战——并且每个任务都通过统一接口暴露，使同一份评估代码评分所有模型。我们以 $\x$ 表示模型输入、$\y$ 表示目标、$\metric$ 表示指标集合。

\paragraph{统一接口。}所有任务实现
\begin{center}\small
\code{split()} $\to$ 时间视图；\quad
\code{featurize(as\_of)} $\to$ 特征；\quad
\code{metrics()} $\to$ 指标列表；\quad
\code{baseline()} $\to$ 注册基线；\quad
\code{evaluate(model)} $\to$ 报告。
\end{center}
\asof{} 参数在特征构造签名中是\emph{强制}的。这是刻意的设计决定：防泄漏是类型层面的义务，而非事后补丁。基线按三个层级组织——平凡、领域与方法论——而方法论层覆盖六大方法家族（统计、树模型、深度、图、贝叶斯与因果），使读者既能看到相对有意义参照点的增益，也能直接比较不同方法家族。

% -----------------------------------------------------------------------------
\subsection{\task{1}：成绩预测（回归）}
\label{sec:task1}

\paragraph{定义。}给定选手赛前的历史成绩序列，预测其在某场比赛、某项目、某轮次中取得的 \code{best} 与 \code{average} 成绩。形式化地，令 $\y = f_\theta(\x)$，其中 $\x = (\x^{\text{seq}}, \x^{\text{static}}, \asof)$，$\y = (\hat{b}, \hat{a}, \hat{\sigma}^2)$ 分别表示预测的最优、平均与预测方差。

\paragraph{输入。}身份字段（选手、比赛、项目、轮次类型）；最近 $N$ 次尝试的历史成绩序列；静态特征（年龄、性别、国籍、参赛总次数）；以及强制传入的决策时刻 \asof{}。

\paragraph{输出。}连续的最优与平均成绩预测（仅当格式支持时才有平均），以及用于区间估计的分布参数。

\paragraph{指标。}因各项目量级差异巨大，我们在\emph{对数域}计算误差：\maelog{}（主指标）与 \rmselog{}（次指标），外加预测区间覆盖率（校准指标）。结果按选手水平分层报告。

\paragraph{基线。}平凡基线：最近 25 次尝试的历史均值。领域/统计基线：近期成绩的线性外推，以及带自助轮次模拟的核密度估计。方法论基线：LSTM/Transformer 等序列模型，以及对数变换目标上的梯度提升回归。

\paragraph{领域挑战。}表现是非平稳的：训练方法改进、伤病与硬件更迭会造成成绩阶跃。项目间方差相差数量级（三阶以秒计、七阶以分钟计、最少步以整数步计），因此项目特定归一化必不可少。

% -----------------------------------------------------------------------------
\subsection{\task{2}：名次预测（排序）}
\label{sec:task2}

\paragraph{定义。}给定某轮次中\emph{全部}参赛选手的历史，预测最终排序。对选手集合 $\mathcal{P}$，学习评分函数 $s_\phi(\x_i)$，输出排列 $\hat{\pi} = \mathrm{argsort}_{i \in \mathcal{P}}(-s_\phi(\x_i))$，以及名次分布 $P(\mathrm{rank}_i = k)$ 与领奖台概率 $P(\mathrm{rank}_i \le 3)$。

\paragraph{输入。}比赛背景（比赛、项目、轮次类型）；每位参赛选手的历史特征；强制传入的 \asof{} 时刻。

\paragraph{输出。}预测的整数名次、名次分布与领奖台概率。

\paragraph{指标。}Kendall $\kendall$ 秩相关（主指标）、与真实领奖台的前 3 重叠度、领奖台概率的 Brier 分数，以及与官方 \wca{} Psych Sheet 的直接对比。

\paragraph{基线。}领域基线：\wca{} \emph{Psych Sheet}，按选手历史最佳 \code{ao5} 排序，忽略方差与稳定性——一个强而实用的参照。统计基线：核密度蒙特卡洛名次模拟与 Plackett--Luce 模型 \citep{plackett_luce}，后者为每位选手分配潜在实力并定义排序分布。方法论基线：在“选手--比赛”异构图上的图神经网络（GNN），建模对抗交互。

\paragraph{领域挑战。}核心难点是交互效应：同一选手面对不同阵容表现可能不同，因此模型必须刻画相对实力而不仅是绝对实力。轮次规则与 \DNF{} 相互作用：\code{ao5} 中单次 \DNF{} 通常作为最差被丢弃，但两次 \DNF{} 会使整轮与名次崩塌。轮次规模从数人到数百人不等，固定长度排序结构难以直接适配。

% -----------------------------------------------------------------------------
\subsection{\task{3}：DNF 预测（不平衡分类）}
\label{sec:task3}

\paragraph{定义。}给定选手在某轮次中已完成的尝试与剩余尝试数，预测其后某次尝试为 \DNF{} 的概率。形式化地，设已观测序列 $\x = (x_1, \dots, x_{k})$、轮次背景 $c$（格式、剩余次数），估计 $P(\text{DNF} \mid \x, c)$ 与整轮 \DNF{} 概率 $P(\text{round DNF})$。

\paragraph{输入。}选手身份；已完成的尝试值（含 \DNF{} 标记）；轮次格式与剩余尝试数；选手历史 \DNF{} 率；强制传入的 \asof{} 时刻。

\paragraph{输出。}每次尝试的 \DNF{} 概率与整轮 \DNF{} 概率。

\paragraph{指标。}\aucroc{} 与 \aucpr{}（不平衡下更强调后者）、F1、Matthews 相关系数（MCC），以及用于校准的可靠性图。由于正类稀有，\aucroc{} 会高估模型能力，我们要求同时报告 \aucpr{}、MCC 及随机/全局率参照。我们额外报告\emph{硬子集}：选手历史 \DNF{} 率位于 $[0.1, 0.3]$，以在基准演化中维持区分度、防止饱和。

\paragraph{基线。}平凡/领域基线：选手历史 \DNF{} 率。统计基线：以已完成尝试均值与方差为特征的 logistic 回归，以及用 Beta 先验做收缩的贝叶斯分层模型。方法论基线：在选手与轮次背景特征上的随机森林/梯度提升。

\paragraph{领域挑战。}\DNF{} 稀有且极不均匀：盲拧与多盲项目失败率远高于三阶，而最少步有其特有的超时/违规失败模式。\DNF{} 的发生还具有序列依赖性：一次失败后选手可能更保守（降低风险）或心态崩溃（提高风险）。最后，风险与规则相关——\code{ao5} 中已有一次 \DNF{} 后，再一次失败的容错空间为零。

% -----------------------------------------------------------------------------
\subsection{\task{4}：人类极限估计（极值外推）}
\label{sec:task4}

\paragraph{定义。}基于各项目历史世界纪录序列 $\{(t_i, y_i)\}$，估计理论表现上限 $L$ 与预期逼近的年份。采用饱和模型
\begin{equation}
  y(t) = L + (y_0 - L)\,\exp\!\big(-\lambda\, g(t)\big) + \varepsilon(t),
  \label{eq:limit}
\end{equation}
其中 $L$ 为极限、$g(t)$ 为技术进步函数、$\varepsilon$ 为噪声，并推断后验 $p(L, \lambda, \text{params} \mid \text{data})$。

\paragraph{输入。}各项目的世界纪录序列（日期、数值），以及项目级特征（均值、方差、历史长度）。

\paragraph{输出。}$L$ 的点估计与后验、预测区间，以及带区间的收敛年份估计。

\paragraph{指标。}极限没有可观测的真值，因此评估依赖：(a) 留一法重拟合下估计的稳定性；(b) 预测区间的经验覆盖率；(c) 与既有领域知识的一致性。作为锚点，速拧社区估计三阶极限约为 $2.37$ 秒、收敛时间约为 2036 年。

\paragraph{基线。}统计基线：式~\ref{eq:limit} 的指数衰减模型。方法论基线：结合极值理论的高斯过程回归 \citep{evt_coles}；将各项目极限视为超分布样本的贝叶斯分层极值模型 \citep{bayesian_hierarchical}；以及识别进步速率结构突变的变点检测。

\paragraph{领域挑战。}首要风险是外推：历史数据无法预见未来由新硬件或新训练体系带来的跃变。数据密度差异极大；三阶有 23 年密集数据，而部分项目的纪录点不足五十个，需要强先验或分层借力。最后，目标是不可观测的，因此验证只能是间接的。

% -----------------------------------------------------------------------------
\subsection{\task{5}：技能迁移（因果推断）}
\label{sec:task5}

\paragraph{定义。}量化某项目表现的提升如何\emph{因果地}影响另一项目的表现。对每个有序项目对 $(A, B)$，估计因果效应 $\theta_{A \to B}$，并组装 $17 \times 17$ 迁移矩阵及置信区间。

\paragraph{输入。}选手的多项目记录（时间戳、项目、成绩），以及每名选手首次参加各项目的时间。

\paragraph{输出。}迁移矩阵 $\Theta = [\theta_{A \to B}]$ 与区间估计，并对无法可靠估计的项目对显式标注“不可识别”。

\paragraph{指标。}相对于留出数据的点估计精度、估计在不同子样本上的稳健性，以及与领域专家判断的一致性。由于缺乏真值效应，我们额外报告\emph{可识别项目对数量}。

\paragraph{基线。}关联基线：项目间成绩的 Pearson/Spearman 相关——忽略混淆的朴素方法。因果基线：将首次参加项目 $A$ 视为处理的双重差分（DID）\citep{did}；以某项目首次在本国举办为工具变量的工具变量法（IV）\citep{iv}；以及用于异质性处理效应的因果森林 \citep{causal_forest}。

\paragraph{领域挑战。}混淆严重：天赋与训练投入会同时提升所有项目表现，造成虚假正迁移。参赛顺序并非随机——擅长项目 $A$ 的选手可能更早接触项目 $B$，从而引入自选择偏差。此外，许多项目对共同参赛选手极少，效应难以识别，因此基准要求显式的“不可识别”标签，而非强行给出估计。
